\chapter{Conclusie}%
\label{ch:conclusie}

\section{Antwoord op de onderzoeksvraag}
\label{sec:antwoord-onderzoeksvraag}

De centrale onderzoeksvraag luidde: \textit{Welke combinatie van chunkingstrategie, embeddingmodel en \gls{llm} resulteert in een optimaal \gls{rag}-systeem voor Turtle S.r.l.'s \gls{poc}?}

Op basis van een systematische grid-search over 27 unieke pipeline-configuraties (3\,×\,3\,×\,3) en 540 evaluatieruns onderscheidt de combinatie \textbf{\texttt{RecursiveCharacterSplitter} + \texttt{multilingual-e5-large-instruct} + \texttt{Gemma-3-27b-it}} zich als de sterkste over alle vijf kwaliteitsmetrics. Deze configuratie is de enige die tegelijkertijd de drempel voor source attribution bereikt (SA\,=\,46{,}7\,\%; $p_{\text{holm}} < 0{,}001$ beter dan alle alternatieven), een sterke context recall behaalt ($\mu_{\text{CR}} = 0{,}792$) en de hoogste faithfulness-trend vertoont ($\mu_F = 0{,}657$). Op basis hiervan wordt de volgende enterprise-architectuur formeel aanbevolen voor de \gls{poc} van Turtle S.r.l.:

\begin{enumerate}
  \item \textbf{Vectorisatie:} \texttt{multilingual-e5-large-instruct} als standaard vector-engine. De meertalige capaciteit, superieure Context Recall ($\mu = 0{,}792$) en de afwezigheid van de SA-penalisering die \texttt{bge-m3} treft, maken het de meest veelzijdige en consistente embedder voor een heterogene, Italiaans- en Engelstalige \gls{esg}-documentstroom. \texttt{Snowflake Arctic Embed} is uitgesloten wegens zijn aantoonbaar lagere retrievalkwaliteit, met direct negatief effect op faithfulness en source attribution stroomafwaarts.
  \item \textbf{Data-ingestie en chunking:} \texttt{RecursiveCharacterSplitter}, wegens zijn statistisch superieure bronvermeldingsscore (46{,}7\,\% correct, $p_{\text{holm}} < 0{,}001$ beter dan de overige strategieën) en zijn structuurbehoudende eigenschap die paginagrensmarkeerders intact laat. In een \gls{csrd}-auditcontext is traceerbaarheid naar exacte brondocumenten prioritair boven maximale context-precisie. Wanneer traceerbaarheid secundair is, biedt de \texttt{FixedSizeTokenSplitter} de hoogste Context Precision en de meest homogene chunk-morfologie.
  \item \textbf{Generatieve audit-engine:} \texttt{Gemma-3-27b-it} op basis van zijn numeriek hoogste faithfulness-gemiddelde ($\mu = 0{,}657$) en source attribution ($\mu = 0{,}344$), ook al zijn deze verschillen statistisch niet significant. In een \gls{esg}-compliancecontext geldt dat bij twijfel de voorkeur gegeven wordt aan het model dat de minste neiging vertoont tot hallucinaties. Het weigeringsgedrag van \texttt{Mistral} is inhoudelijk integer maar drukt de SA-score structureel wanneer de upstream pipeline niet optimaal presteert; \texttt{Llama} vertoont de zwakste faithfulness-trend ($\mu = 0{,}583$).
\end{enumerate}

Het meest opmerkelijke empirische inzicht van dit onderzoek is de hiërarchische structuur van componenteffecten: het \textbf{embeddingmodel} is de dominante kwaliteitshefboom, de \textbf{chunkingstrategie} is doorslaggevend voor bronvermelding, maar het \textbf{\gls{llm}} vertoont geen statistisch significant effect op faithfulness noch op answer relevancy ($H = 5{,}78$, $p = 0{,}056$ resp.\ $H = 5{,}00$, $p = 0{,}082$). Dit is een fundamenteel inzicht voor de implementatiepraktijk: optimalisatie-inspanningen bij \gls{rag}-pipelines voor \gls{esg}-compliance zijn beter gericht op de upstream vectorisatie- en chunkfase dan op de selectie van een krachtigere generatieve engine.

Een tweede structureel gegeven is de incompatibiliteit tussen hoge source attribution en hoge answer relevancy: de enige embedder die AR\,$\geq$\,0{,}65 haalt (\texttt{bge-m3}, $\mu = 0{,}679$) vertoont tegelijk de laagste SA ($\mu = 0{,}233$). Omgekeerd vereist SA\,$\geq$\,40\,\% \texttt{E5-large} als embedder, wat AR tot $\mu = 0{,}635$ beperkt. Deze trade-off is geen artefact van hyperparameterafstemming maar een structurele eigenschap van de hybride multi-vector RRF-retrieval van \texttt{bge-m3}: de brede semantische dekking die de antwoordrelevantie verhoogt, verdunt tegelijk de paginaspecifieke metadata-koppeling die source attribution vereist.

Wetenschappelijk levert dit onderzoek de eerste systematische componentvergelijking op van een \gls{rag}-pipeline die source attribution als primaire auditmetriek behandelt in een meertalige \gls{esg}/\gls{csrd}-context. De bevinding dat het \gls{llm}-effect statistisch onbeduidend is ten opzichte van de upstream componenten nuanceert de gangbare aanname dat generatieve modelverbetering de voornaamste kwaliteitshefboom is, en verschuift het optimalisatiediscours naar de vectorisatie- en chunkfase.

\section{Kritische reflectie}
\label{sec:kritische-reflectie}

De uitkomsten van dit onderzoek liggen deels in lijn met de verwachtingen uit de literatuurstudie, maar bevatten ook meerdere verrassende bevindingen die de gangbare aannames over \gls{rag}-optimalisatie nuanceren.

\paragraph{Het \glstext{llm} als minst bepalende factor.}
Het meest onverwachte resultaat is het ontbreken van een statistisch significant \gls{llm}-effect. In de literatuur en praktijk wordt de \gls{llm}-keuze doorgaans gepresenteerd als de primaire kwaliteitsdriver van een \gls{rag}-systeem. Dit onderzoek falsifieert die aanname voor de specifieke context van \gls{esg}-compliance-auditing: bij voldoende contextomvang (top-$k = 5$ chunks) worden de generatieve kwaliteitsverschillen tussen Gemma, Mistral en Llama statistisch geneutraliseerd. Dit benadrukt dat het praktische optimalisatierendement hogerop in de ketting ligt — in de vectorisatie- en chunkfase — en niet in de voortdurende vervanging van het \gls{llm} door krachtigere varianten.

\paragraph{De cascade-failure van Snowflake.}
De zwakke prestaties van \texttt{Snowflake Arctic Embed} waren deels te verwachten op basis van zijn lagere trainingsdiversiteit, maar de omvang van het cascade-effect was opmerkelijk. Slechte retrieval door \texttt{Snowflake} triggerde de Mistral-weigeringsparadox — waarbij het \gls{llm} correct weigert een antwoord te genereren, maar de geautomatiseerde SA-check dit registreert als een SA-score van nul — en drukte zo de evaluatieresultaten van de volledige Snowflake-Mistral-combinatie disproportioneel. Dit illustreert hoe een zwakke upstream component onverwachte side-effects kan induceren bij specifieke LLM-gedragsprofielen.

\paragraph{Betrouwbaarheid van de evaluatiemetrics.}
De menselijke validatiestudie (Sectie~\ref{sec:methodologische-validatie}) toont dat de vijf \gls{ragas}-metrics een sterk variërend betrouwbaarheidsniveau hebben. \textit{Context Recall} vertoonde een uitzonderlijk sterke mens-machine-uitlijning ($\rho = 0{,}973$), waardoor vergelijkingen op basis van CR robuust zijn. \textit{Faithfulness} en \textit{Answer Relevancy} vertonen daarentegen zwakke correlaties ($\rho = 0{,}257$ resp.\ $\rho = 0{,}460$) met de menselijke beoordeling, mede door de taalkundige complexiteit van Italiaanse technisch-juridische terminologie. Dit impliceert dat absolute drempeloordelen op basis van F en AR met gepaste voorzichtigheid moeten worden geïnterpreteerd, en dat de relatieve rangvergelijkingen tussen configuraties betrouwbaarder zijn dan de absolute scores.

\paragraph{De praktische bijdrage.}
De beslisboom in Sectie~\ref{subsec:beslisboom} vormt de centrale praktische bijdrage van dit onderzoek. In tegenstelling tot een generieke aanbeveling op basis van één rangschikking, laat de boom toe dat toekomstige implementatoren de configuratie afstemmen op hun specifieke inzetcontext — afhankelijk van of juridische traceerbaarheid (SA), antwoordrelevantie (AR) of volledigheid van de bewijslast (CR) de primaire eis is. Dit verplaatst de optimalisatiestrategie van een éénmalige \glspl{llm}-selectiebeslissing naar een iteratieve, metriekgestuurde componentkeuze.

\paragraph{Kosten en latency bewust buiten scope.}
Latency en operationele kosten zijn in dit onderzoek niet gemeten en bewust buiten de evaluatiescope gehouden. In een \gls{csrd}-auditcontext — waar elke gegenereerde bewering wettelijk traceerbaar moet zijn — mag kwaliteit niet worden opgeofferd voor snelheid of kostenreductie. Een lagere faithfulness of source attribution omwille van snellere inference is in een compliancecontext geen aanvaardbare trade-off. Latency-optimalisatie is bijgevolg een secundaire operationele zorg die pas relevant wordt nadat de kwaliteitsdrempels zijn gehaald, en vormt een afzonderlijk engineeringtraject.

\paragraph{LLM-gebaseerde chunkingstrategieën uitgesloten.}
Agentic en LLM-gebaseerde chunkingstrategieën — zoals propositionele splitsing waarbij een \gls{llm} de documentinhoud herschrijft tot atomaire beweringen — werden bewust uitgesloten uit de shortlist. De voornaamste reden is methodologisch: wanneer het \gls{llm} hallucinaties introduceert tijdens de chunkingfase, worden deze fouten als grondwaarheid opgeslagen in de vector database. In die situatie is de hallucinatiegevoeligheid niet langer geïsoleerd tot de generatieve fase, maar verspreid over twee componenten van de pipeline. Dit maakt het statistisch onmogelijk om het effect van de \gls{llm}-component in de generatieve fase te isoleren, wat de vergelijkbaarheid van configuraties zou ondermijnen.

\section{Beperkingen en verder onderzoek}
\label{sec:beperkingen}

Dit onderzoek kent een aantal methodologische beperkingen die de generaliseerbaarheid van de resultaten inperken en aanleiding geven tot nader onderzoek.

\paragraph{Corpusbeperkingen.}
De experimentele documentbasis bestond uitsluitend uit de interne duurzaamheidsdocumenten van Turtle S.r.l., opgesteld in het Italiaans en Engels. De resultaten zijn bijgevolg direct toepasbaar voor vergelijkbare heterogene, meertalige \gls{esg}-documentstromen, maar de generaliseerbaarheid naar andere domeinen (juridische contracten, medische dossiers, technische handleidingen) vereist afzonderlijke validatie. Onderzoek met een grotere en meer diverse documentbasis zou de externe validiteit van de geobserveerde componenthiërarchie kunnen bevestigen of bijstellen.

\paragraph{Omvang van de gouden dataset.}
De evaluatiedataset telde 20 handmatig geverifieerde vraag-antwoordparen, wat voor de statistische vergelijkingen voldoende was maar de diversiteit van mogelijke \gls{esg}-vraagtypen slechts gedeeltelijk dekt. Een uitbreiding naar een rijkere golden dataset — met vragen over impliciete verbanden, kwantitatieve indicatoren en cross-documentreferenties — zou genuanceerdere uitspraken mogelijk maken over de generaliseerbaarheid van de resultaten.

\paragraph{Vaste hyperparameters.}
De experimentele configuraties varieerden uitsluitend de componentkeuze (chunker, embedder, \gls{llm}); hyperparameters zoals chunkgrootte (1024 tokens), overlapfactor (154 tokens) en de top-$k$-retrievaldiepte (3 chunks) werden gefixeerd. Het is aannemelijk dat het optimale componentprofiel verschuift bij andere hyperparameterinstellingen — met name bij grotere chunks of een hogere $k$ — waardoor de huidige drempelwaarden als richtinggevend en niet als universeel gelden.

\paragraph{Preprocessingpipeline en paginanummering.}
De preprocessing-pipeline was gefocust op de extractie van platte tekst en het injecteren van paginagrensmarkeerders (\texttt{<<<PAGE X>>>}). Bij Word-documenten (\texttt{.docx}) introduceerde de conversiestap een vaste offset van twee à drie pagina's ten opzichte van de werkelijke paginanummering, doordat kopteksten en inhoudsopgaven als extra paginaoppervlak worden meegerekend. Dit leidt tot een systematische onderschatting van de source attribution-scores voor de documenten in Word-formaat, omdat de gerapporteerde paginanummers niet overeenkomen met de nummering in het originele bestand. Een robuustere preprocessing-pipeline die uniform de fysieke paginanummering extraheerde, zou dit artefact elimineren maar viel buiten de scope van dit onderzoek.

\paragraph{Juridische en privacybeperkingen (Cloud Act).}
De \gls{poc} van Turtle S.r.l. zal worden gehost op Microsoft Azure, een Amerikaans cloudplatform dat onder de jurisdictie van de \textit{Clarifying Lawful Overseas Use of Data Act} (Cloud Act) valt. Dit betekent dat de Amerikaanse overheid in principe toegang kan eisen tot bedrijfsdata die op Azure-servers zijn opgeslagen, ongeacht of die servers zich fysiek in de \gls{eu} bevinden. Voor \gls{esg}-documenten die vertrouwelijke bedrijfsinformatie bevatten, stelt dit een potentieel conflict met de \gls{gdpr} en de Europese datavereisten voor \gls{csrd}-compliance. Een mogelijke mitigatie is een anonymiseringsstap in de preprocessing-pipeline die persoonsgegevens en commercieel gevoelige informatie maskeert vóór ingestie in de vector database. De ontwikkeling en validatie van een dergelijke anonimiseringsmodule valt echter buiten de scope van dit onderzoek en vormt een relevant vervolgtraject.

\paragraph{Toekomstig onderzoek.}
Op basis van de bevindingen worden vier onderzoeksrichtingen als meest veelbelovend beschouwd:

\begin{enumerate}
  \item \textbf{Domeinspecifieke fine-tuning van embedders.} De meertalige capaciteit van \texttt{E5-large} en de hybride retrieval van \texttt{bge-m3} zijn generieke eigenschappen. Fine-tuning op een corpus van \gls{esrs}-documenten, \gls{csrd}-rapportagestandaarden en juridische \gls{esg}-teksten zou de context recall en source attribution verder kunnen verbeteren door de representatieruimte te aligneren met domeinspecifieke terminologie.
  \item \textbf{Multi-hop reasoning voor cross-documentvragen.} De huidige single-hop retrieval (één queryvector → top-$k$ chunks) is ontoereikend voor complexe \gls{esg}-vragen waarbij bewijs verspreid is over meerdere documenten of impliciete redeneerketens vereisen. \gls{rag}-architecturen met iteratieve of gestructureerde multi-hop retrieval bieden een logische uitbreiding voor diepgaandere compliance-analyse.
  \item \textbf{Productievalidatie en gebruikersstudie.} De benchmark-resultaten zijn gegenereerd onder gecontroleerde laboratoriumomstandigheden. Een longitudinale validatiestudie waarbij auditprofessionals de gegenereerde antwoorden beoordelen in een werkelijke \gls{csrd}-rapportagecontext, zou inzicht geven in de operationele bruikbaarheid en de discrepantie tussen geautomatiseerde \gls{ragas}-scores en professioneel oordeel.
  \item \textbf{Robuuste preprocessingpipeline met privacybescherming.} Een geïntegreerde preprocessingmodule die paginanummering correct normaliseert over bestandsformaten en tegelijk een anonimiseringsstap uitvoert, zou zowel de SA-nauwkeurigheid verhogen als de Cloud Act-compliance-problematiek adresseren. Dit vormt een noodzakelijk engineeringtraject voor de productietransitie van de \gls{poc}.
\end{enumerate}
